\documentclass[11pt,a4paper]{article}
\usepackage[margin=23mm]{geometry}
\usepackage[T1]{fontenc}
\usepackage{lmodern}
\usepackage{graphicx}
\newcommand{\evidence}{../results/relaxed-last-switch/error-scale-3x}
\setlength{\parindent}{0pt}

\begin{document}
\section*{Relaxed final-boundary comparison}

All three methods use the same gain and threefold initial error. The panels compare anchor and fixed-clock errors, individual cycle durations, and first-cycle dwells. Blue denotes \texttt{daniel}, green \texttt{daniel-relaxed}, and orange \texttt{marcolino}.

\begin{center}
\includegraphics[width=\linewidth,height=.78\textheight,keepaspectratio]{\evidence/comparison.pdf}
\end{center}

\clearpage
The first 25 cycles show error convergence, conditioning activity, and requested versus applied minimum dwells. The two Daniel policies share a beta axis; Marcolino uses a separate delayed-switch count. All applied dwells meet the $3\,\mu\mathrm{s}$ bound.

\begin{center}
\includegraphics[width=\linewidth,height=.86\textheight,keepaspectratio]{\evidence/article_response_with_marcolino.pdf}
\end{center}

\clearpage
Exact state trajectories over the same $3.432$ ms window compare all three methods with the dashed fixed-clock reference. The panels show the state-space paths, both capacitor voltages, and inductor current; each method retains its own sampling times.

\begin{center}
\includegraphics[width=\linewidth,height=.86\textheight,keepaspectratio]{\evidence/state_trajectories.pdf}
\end{center}
\end{document}
